\BeleskeID{01-s032}
% element: 01-s032:e01
% element: 01-s032:e02

% element: 01-s032:d01
Za EXILI neparni broj jedinica označen je izlazom 1, a za EXNILI izlazom 0. Pri čitanju reči „aktivan“ zato najpre treba odrediti da li je aktivan nivo jedinica ili nula.

U donjoj desnoj realizaciji prvi EXILI formira paritet prva dva ulaza, \(p=x_1\oplus x_2\), a izlazni EXNILI daje
\[
y=\overline{p\oplus x_3}
  =\overline{x_1\oplus x_2\oplus x_3}.
\]
Negacija se ovde uvodi jednom, na kraju, pa izlaz odgovara direktnom troulaznom EXNILI kolu. Za broj ulaznih jedinica 0, 1, 2 i 3 izlazi EXILI su redom 0, 1, 0 i 1, dok su izlazi EXNILI redom 1, 0, 1 i 0. Time se mogu proveriti oba para šema.
